% ==================================================================================================================================
% Three pieces of text for hires2_correction.tex, on the arrangement of the material along the line of sight.
%
%  1. A replacement for the last paragraph of Sect. 5.3 (\label{sec:orient}), the one beginning "One assumption remains".
%  2. A new subsection of Sect. 8 (\label{sec:limits}), to follow "The geometry".
%  3. A new appendix, to follow "The grid of models".
%
% All the numbers are produced by superposition_test.py, which writes superposition_test.txt.
% ==================================================================================================================================


% ----------------------------------------------------------------------------------------------------------------------------------
% 1. REPLACES the last paragraph of Sect. 5.3, "One assumption remains, and it is not small. ... Section~\ref{sec:limits} returns
%    to this."
% ----------------------------------------------------------------------------------------------------------------------------------

One assumption remains, that each structure is of uniform density, and it matters far less than it appears. Along a
plane-parallel line of sight it does not matter at all. What shields a parcel of dust is the column of material lying between it
and the surface, and that column is what it is however the material within it is distributed: piling half of it into a dense clump
changes nothing about how much of it there is. The mass of a line of sight per unit column is constant as well, for any density
profile whatever, because the column is itself the integral of the density. Writing $\Sig(s)$ for the column in front of a parcel
at the depth $s$, the emission of Eq.~(\ref{eq:emission}) integrated along the line of sight is
\begin{equation}
   \int_0^L B_\lambda\big(T(\Sigma_{\rm a}(\Sig(s)))\big)\,\rho(s)\,{\rm d}s
   = \int_0^{\Sig} B_\lambda\big(T(\Sigma_{\rm a}(\Sig'))\big)\,{\rm d}\Sig' ,
\label{eq:invariance}
\end{equation}
and the density has disappeared. Two structures of the same total column emit the same spectrum whether the material is smooth or
clumped, and whether a clump sits at the near face, in the middle or at the far face. Appendix~\ref{app:arrange} confirms this
numerically on profiles differing by a factor of thirty in density, which reproduce the smooth case to two parts in
$10^{4}$.

Concentration and arrangement can therefore act only through the one thing Eq.~(\ref{eq:invariance}) assumes, that the radiation
leaves along the line of sight rather than sideways. They do, and by a modest amount. A sphere whose central-to-edge density ratio
is 10 requires a correction 3 per cent larger than a uniform sphere of the same central column, and one whose ratio is $10^3$
requires between 6 and 15 per cent more, the difference growing with the column. A cirrus with a denser structure embedded in it
requires between 15 per cent less and 11 per cent more than the single uniform body that the method substitutes for it, depending
on where within the cirrus the structure lies. Both are quantified in Appendix~\ref{app:arrange} and enter the budget of
Sect.~\ref{sec:limits}.


% ----------------------------------------------------------------------------------------------------------------------------------
% 2. NEW SUBSECTION of Sect. 8, to follow "The geometry".
% ----------------------------------------------------------------------------------------------------------------------------------

\subsection{The arrangement along the line of sight}

An observed line of sight is not one structure but several, at depths that no observation reveals, and the method replaces them by
a single uniform body of the same total column. Equation~(\ref{eq:invariance}) shows that this substitution is exact in the
plane-parallel limit, so the error it commits is confined to the transverse escape of bounded structures.
Appendix~\ref{app:arrange} measures it in the two forms it takes. Concentrating a bounded structure raises the required correction
by 3 per cent at a central-to-edge density ratio of 10 and by up to 15 per cent at a ratio of $10^3$. Embedding a structure in a
cirrus changes it by between $-15$ and $+11$ per cent according to the depth at which the structure sits, without a preferred
sign, and the value the method applies falls inside that range at every column we examined. We take 15 per cent as the
uncertainty from this source. It is comparable to the uncertainty of the shielding relation and larger than that of the geometry,
and unlike either of them it cannot be reduced by a better calculation, because the depth of a structure within a cloud is not an
observable.


% ----------------------------------------------------------------------------------------------------------------------------------
% 3. NEW APPENDIX, to follow "The grid of models".
% ----------------------------------------------------------------------------------------------------------------------------------

\section{The arrangement of the material along the line of sight}
\label{app:arrange}

The correction is applied pixel by pixel, and a pixel receives the emission of everything along its line of sight. A single
modified blackbody is fitted to that sum, so what biases the fit is the whole distribution of dust temperatures along the line of
sight, and not the temperature of any one structure within it. The method replaces that distribution by the one that a single
uniform body of the same total column would have. This appendix measures what the replacement costs. All of it is computed by
\texttt{superposition\_test.py}, which takes the relation of Eq.~(\ref{eq:law}), the angular average of Eq.~(\ref{eq:seff}) and
the constants from the implementation itself, so that nothing is duplicated.

\subsection{The invariance}

For a plane-parallel geometry the emergent emission depends on the total column alone. The attenuating column of a parcel is
determined by the column between it and each surface, which does not depend on how that column is distributed, and the mass per
unit column of a line of sight is constant for any density profile, which is Eq.~(\ref{eq:invariance}). As a check on the algebra
and on the integrations, we computed the four intensities of a line of sight of $2{\times}10^{22}$\,cm$^{-2}$ carrying a clump
thirty times the ambient density, placed at 15, 50 and 85 per cent of the depth. They differ from the smooth case of the same
total column by $2.1{\times}10^{-4}$, $4.6{\times}10^{-5}$ and $1.2{\times}10^{-4}$ at most, which is the accuracy of the
quadrature. Neither the concentration of the material nor the position of the clump has any effect.

\subsection{Concentration}

Since the density profile leaves the problem entirely in the plane-parallel limit, concentration can act only through the escape
of radiation across the line of sight, which a bounded structure permits and a layer does not. We measure it on a sphere of
density $\rho(r) \propto (1 + r^2/a^2)^{-1}$ out to its boundary, normalized so that the column through its center is the observed
one, with $a$ set by the central-to-edge density ratio. The parcels of the central line of sight are weighted by their density,
which is how the emission of that line of sight weights them, and the column each of them must send its radiation through is
integrated along the ray to the boundary in every direction of the average.

\begin{table}
\caption{Correction factor required by the central line of sight of a sphere of the given central column, for four
central-to-edge density ratios. The first column is the uniform sphere that the method assumes.}
\label{tab:concentr}
\centering
\begin{tabular}{lcccc}
\hline\hline
\noalign{\smallskip}
$\Sig$ (cm$^{-2}$) & uniform & $10$ & $10^2$ & $10^3$ \\
\noalign{\smallskip}
\hline
\noalign{\smallskip}
$5{\times}10^{21}$ & 1.068 & 1.096 & 1.122 & 1.136 \\
$1{\times}10^{22}$ & 1.218 & 1.260 & 1.316 & 1.351 \\
$3{\times}10^{22}$ & 1.709 & 1.750 & 1.857 & 1.968 \\
$1{\times}10^{23}$ & 2.288 & 2.327 & 2.438 & 2.635 \\
\noalign{\smallskip}
\hline
\end{tabular}
\end{table}

Table~\ref{tab:concentr} gives the result. Concentration always raises the required correction, because the mass of a
concentrated structure sits where the shielding is deepest, and the effect grows with the column, from 3 per cent at a ratio of 10
to 15 per cent at a ratio of $10^3$. It is bounded and one-sided, and a correction built on uniform structures therefore
underestimates the deficit of a concentrated one, never the reverse.

\subsection{Two superposed components}

The second form of the question is a line of sight crossing a diffuse cirrus and a denser structure embedded in it at an unknown
depth. We take a plane-parallel cirrus of column $\Sig_{\rm c}$, in which the material at the column depth $x$ is shielded by the
angular average of $x/|\mu|$ toward the observer and $(\Sig_{\rm c} - x)/|\mu|$ away from it, and a uniform sphere whose column
through its center is $\Sig_{\rm s}$, embedded in that cirrus with the fraction $q$ of the cirrus column in front of it. A parcel
at the signed position $z$ along the diameter of the sphere, in units of its radius, is shielded by its own sphere over the path
\begin{equation}
   l(z,\psi) = -z\cos\psi + \sqrt{1 - z^2\sin^2\psi} ,
\label{eq:lsphere}
\end{equation}
in units of the radius, and by the cirrus over $q\,\Sig_{\rm c}/|\cos\psi|$ toward the observer and
$(1-q)\,\Sig_{\rm c}/|\cos\psi|$ away from it. The sign of $z$ has to be kept: with $|z|$ in its place the sphere would be
symmetric about its center whatever lay in front of it, and the answer would be wrong for every $q$ but one half. The angular
average of Eq.~(\ref{eq:seff}) is then taken over $\psi$ with the isotropic weight $\sin\psi$, the temperature of every parcel
follows from Eq.~(\ref{eq:law}), the emission of the cirrus and of the sphere are added with their masses, and one modified
blackbody is fitted to the sum exactly as the reconstruction fits it. The correction factor is the ratio of the true column to the
fitted one, as in Eq.~(\ref{eq:factor}). The construction is symmetric under $q \rightarrow 1-q$, as it must be, and we have
verified that the computed factors are.

\begin{table}
\caption{Correction factor required by a cirrus of column $\Sig_{\rm c}$ with a sphere embedded in it carrying the rest of the
column, against the fraction $q$ of the cirrus lying in front of the sphere, beside the factor the method applies to the same
total column. The values at $q$ and at $1-q$ are equal, so only half the range is listed.}
\label{tab:arrange}
\centering
\begin{tabular}{llcccc}
\hline\hline
\noalign{\smallskip}
$\Sig$ (cm$^{-2}$) & $\Sig_{\rm c}$ (cm$^{-2}$) & method & $q\,{=}\,0$ & $q\,{=}\,0.25$ & $q\,{=}\,0.5$ \\
\noalign{\smallskip}
\hline
\noalign{\smallskip}
$1{\times}10^{22}$ & $1{\times}10^{21}$ & 1.266 & 1.297 & 1.308 & 1.309 \\
                   & $3{\times}10^{21}$ & 1.266 & 1.170 & 1.253 & 1.280 \\
$2{\times}10^{22}$ & $1{\times}10^{21}$ & 1.538 & 1.643 & 1.612 & 1.605 \\
                   & $3{\times}10^{21}$ & 1.538 & 1.538 & 1.541 & 1.547 \\
                   & $1{\times}10^{22}$ & 1.538 & 1.300 & 1.486 & 1.557 \\
$5{\times}10^{22}$ & $1{\times}10^{21}$ & 2.004 & 2.206 & 2.142 & 2.127 \\
                   & $3{\times}10^{21}$ & 2.004 & 2.160 & 2.044 & 2.027 \\
                   & $1{\times}10^{22}$ & 2.004 & 1.945 & 1.967 & 2.001 \\
$1{\times}10^{23}$ & $1{\times}10^{21}$ & 2.294 & 2.548 & 2.475 & 2.453 \\
                   & $3{\times}10^{21}$ & 2.294 & 2.548 & 2.377 & 2.344 \\
                   & $1{\times}10^{22}$ & 2.294 & 2.447 & 2.286 & 2.294 \\
\noalign{\smallskip}
\hline
\end{tabular}
\end{table}

Table~\ref{tab:arrange} gives the result and Fig.~\ref{fig:superposition} shows where it comes from. The departure from the value
the method applies runs from $-15$ to $+11$ per cent, it has no preferred sign, and the method's value lies inside the range
spanned by the arrangement at every column. The spread is largest, 20 per cent, when the diffuse component carries half of the
column, and falls below 10 per cent in most of the configurations we examined.

The middle panel of Fig.~\ref{fig:superposition} shows why the effect exists at all, and that it is a property of the fit rather
than of the geometry. The two lines of sight drawn there carry exactly the same total column of $2{\times}10^{22}$\,cm$^{-2}$, but
in one of them half of it is a sphere pressed against the near face of the cirrus, where it loses radiation sideways and has
nothing in front of it to block the incoming field. That arrangement is warmer, it is brighter in all four wavebands, and the
single modified blackbody fitted to it returns $1.54{\times}10^{22}$\,cm$^{-2}$ where the uniform body returns
$1.30{\times}10^{22}$\,cm$^{-2}$. The same true column would then need a correction of 1.30 rather than 1.54.

Nothing observable distinguishes the two. The depth of a structure within a cloud is not a measurable quantity, and no
improvement of the calculation can recover it, so the 15 per cent of Sect.~\ref{sec:limits} is an irreducible uncertainty of any
method of this kind rather than a defect of this one.
